\chapter{Holonomia e o Teorema de Ambrose--Singer}

\section{Holonomia: Defini\c{c}\~ao Formal}

O conceito de holonomia captura, em linguagem geom\'etrica, a ideia fundamental de que o transporte paralelo ao longo de um ciclo fechado nem sempre devolve o objeto transportado ao seu estado original. Em Teoria Geom\'etrica de Arbitragem, esse desvio constitui precisamente o ganho (ou perda) l\'iquido de uma estrat\'egia de trading que retorna \`a mesma configura\c{c}\~ao de mercado.

\subsection{Conex\~ao e Transporte Paralelo}

Seja \(P(M,G)\) um fibrado principal sobre uma variedade base \(M\) (o espa\c{c}o de estados do mercado) com grupo estrutural \(G\) (o grupo de gauge). Denotamos por \(\chi\) uma conex\~ao sobre \(P\), i.e., uma 1-forma em \(P\) com valores na \'algebra de Lie \(\mathfrak{g}\) satisfazendo:

\begin{enumerate}
    \item \(\chi(A^*) = A\) para todo \(A \in \mathfrak{g}\), onde \(A^*\) \'e o campo fundamental associado.
    \item \(R_g^* \chi = \operatorname{Ad}_{g^{-1}} \chi\) para todo \(g \in G\).
\end{enumerate}

Dada uma curva \(\gamma: [0,1] \to M\) e um ponto \(u \in P_{\gamma(0)}\) na fibra sobre \(\gamma(0)\), o levantamento horizontal de \(\gamma\) com condi\c{c}\~ao inicial \(u\) \'e a \'unica curva \(\tilde{\gamma}: [0,1] \to P\) tal que \(\pi \circ \tilde{\gamma} = \gamma\), \(\tilde{\gamma}(0) = u\), e \(\chi(\dot{\tilde{\gamma}}(t)) = 0\) para todo \(t\).

\subsection{Grupo de Holonomia}

\begin{defi}[Grupo de Holonomia]
Seja \(u \in P\) um ponto fixado. O \textbf{grupo de holonomia} de \(\chi\) com ponto base \(u\) \'e o subgrupo de \(G\) definido por
\begin{equation}\label{eq:hol-group}
    \operatorname{Hol}_u(\chi)
    = \{\, g \in G \mid \exists\; \gamma \text{ la\c{c}o em } M \text{ baseado em } \pi(u) \text{ tal que } \tilde{\gamma}(1) = u \cdot g \,\}.
\end{equation}
\end{defi}

Em palavras: \(\operatorname{Hol}_u(\chi)\) consiste de todos os elementos de \(G\) que podem ser realizados como o desvio de paralelismo ao transportar horizontalmente \(u\) ao longo de algum la\c{c}o fechado na base \(M\). A mudan\c{c}a de ponto base segue a regra de conjuga\c{c}\~ao: \(\operatorname{Hol}_{u \cdot g}(\chi) = g^{-1} \operatorname{Hol}_u(\chi) \, g\). Por essa raz\~ao, frequentemente escrevemos \(\operatorname{Hol}(\chi)\) sem refer\^encia expl\'icita ao ponto base, entendendo o grupo a menos de isomorfismo.

\begin{defi}[Holonomia Local e Restrita]
\begin{enumerate}
    \item A \textbf{holonomia restrita} \(\operatorname{Hol}_u^{\,0}(\chi)\) \'e o subgrupo obtido considerando apenas la\c{c}os \textbf{contr\'ateis} (homot\'opicos a um ponto). Este \'e exatamente a componente conexa da identidade de \(\operatorname{Hol}_u(\chi)\).
    \item A \textbf{holonomia local} \'e a holonomia restrita quando restringimos a aten\c{c}\~ao a vizinhan\c{c}as suficientemente pequenas de um ponto \(x \in M\).
\end{enumerate}
\end{defi}

\begin{prop}
\(\operatorname{Hol}_u^{\,0}(\chi)\) \'e um subgrupo de Lie conexo de \(G\). O quociente \(\operatorname{Hol}_u(\chi) / \operatorname{Hol}_u^{\,0}(\chi)\) \'e discreto e isomorfo a um subgrupo do grupo fundamental \(\pi_1(M)\).
\end{prop}

\begin{proof}[Esbo\c{c}o]
A demonstra\c{c}\~ao cl\'assica (Kobayashi--Nomizu, 1963, Cap\'itulo II, Teorema 4.2) constr\'oi uma subvariedade de \(P\) chamada \textbf{fibrado de holonomia} \(P(u) \subset P\), cujas fibras s\~ao precisamente as \'orbitas de \(\operatorname{Hol}_u(\chi)\). A redu\c{c}\~ao do fibrado principal a \(P(u)\) fornece a estrutura de subgrupo de Lie. A rela\c{c}\~ao com \(\pi_1(M)\) vem da aplica\c{c}\~ao natural que associa a cada la\c{c}o o elemento de holonomia correspondente, cujo n\'ucleo cont\'em os la\c{c}os com holonomia trivial.
\end{proof}

\subsection{Intui\c{c}\~ao Financeira}

A analogia financeira \'e direta e esclarecedora. Considere o seguinte cen\'ario: um trader parte de uma configura\c{c}\~ao de portf\'olio inicial, executa uma sequ\^encia de opera\c{c}\~oes de compra e venda que, ao final, restaura exatamente os mesmos pre\c{c}os de mercado do in\'icio (um la\c{c}o fechado no espa\c{c}o de estados). Se a conex\~ao \(\chi\) (a estrutura de precifica\c{c}\~ao) possui holonomia n\~ao trivial, o portf\'olio final n\~ao coincide com o inicial -- a diferen\c{c}a \'e precisamente o lucro (ou preju\'izo) de arbitragem. Em suma:

\[
\boxed{\text{Ciclo fechado de trading} \;\longleftrightarrow\; \text{La\c{c}o em } M}
\qquad
\boxed{\text{Ganho l\'iquido} \;\longleftrightarrow\; \text{Holonomia n\~ao trivial}}
\]

\section{Parametriza\c{c}\~ao de Estrat\'egias pela Holonomia}

\subsection{Correspond\^encia Fundamental}

A rela\c{c}\~ao central da Teoria Geom\'etrica de Arbitragem estabelece uma correspond\^encia bijetiva entre estrat\'egias de arbitragem e a \'algebra de Lie do grupo de holonomia.

\begin{prop}[Parametriza\c{c}\~ao de Estrat\'egias]
Seja \(\mathfrak{hol}(\chi) = \operatorname{Lie}(\operatorname{Hol}(\chi))\) a \'algebra de Lie do grupo de holonomia. Cada elemento \(A \in \mathfrak{hol}(\chi)\) corresponde a uma classe de equival\^encia de estrat\'egias de arbitragem, onde duas estrat\'egias s\~ao equivalentes se produzem o mesmo ganho infinitesimal em primeira ordem.
\end{prop}

\subsection{Por que Funciona}

A justificativa se desenvolve em tr\^es etapas conceituais:

\begin{enumerate}
    \item \textbf{Transporte ao longo de \(\gamma \to\) elemento do grupo.} Dado um la\c{c}o \(\gamma\) em \(M\) baseado em \(x_0\), o transporte paralelo ao longo de \(\gamma\) define um isomorfismo da fibra \(P_{x_0}\) sobre si mesma, i.e., um elemento \(g_\gamma \in \operatorname{Hol}_{u}(\chi) \subset G\). Este elemento codifica o efeito acumulado da estrat\'egia de trading que realiza o ciclo \(\gamma\).

    \item \textbf{Curvas homot\'opicas \(\to\) mesma holonomia restrita.} Se dois la\c{c}os \(\gamma_0\) e \(\gamma_1\) s\~ao homot\'opicos (com extremos fixos), seus levantamentos horizontais terminam no mesmo ponto da fibra. Logo, pertencem \`a mesma classe em \(\operatorname{Hol}_u^{\,0}(\chi)\). Isto significa que estrat\'egias de trading que podem ser deformadas continuamente uma na outra -- sem ``atravessar singularidades'' do mercado -- produzem o mesmo resultado financeiro.

    \item \textbf{\'Algebra de Lie \(\leftrightarrow\) classes infinitesimais.} Elementos da \'algebra de Lie \(\mathfrak{hol}(\chi)\) s\~ao geradores infinitesimais de elementos do grupo. No contexto financeiro, correspondem a estrat\'egias de arbitragem ``instant\^aneas'' ou de primeira ordem, a partir das quais estrat\'egias finitas podem ser reconstru\'idas por integra\c{c}\~ao (exponencia\c{c}\~ao na \'algebra de Lie).
\end{enumerate}

\begin{prop}[Invari\^ancia sob Homotopia]
Se \(\gamma_0 \simeq \gamma_1\) (homotopia com extremos fixos), ent\~ao \(\tilde{\gamma}_0(1) = \tilde{\gamma}_1(1)\) e, portanto, \(g_{\gamma_0} = g_{\gamma_1}\) em \(\operatorname{Hol}_u^{\,0}(\chi)\).
\end{prop}

\begin{proof}
Decorre diretamente do teorema de unicidade do levantamento horizontal: a homotopia \(H: [0,1] \times [0,1] \to M\) entre \(\gamma_0\) e \(\gamma_1\) levanta-se a uma homotopia horizontal \(\tilde{H}\) em \(P\) com a mesma condi\c{c}\~ao inicial, e a avalia\c{c}\~ao em \(t=1\) \'e constante ao longo do par\^ametro de homotopia \(s\) (Kobayashi--Nomizu, 1963, Cap. II, Prop.~3.2).
\end{proof}

\section{O Teorema de Ambrose--Singer}

\subsection{Enunciado}

O Teorema de Ambrose--Singer (Ambrose \& Singer, 1953) \'e o resultado estrutural mais importante para a Teoria Geom\'etrica de Arbitragem, pois conecta diretamente a curvatura (que mede a ``n\~ao integrabilidade'' local da distribui\c{c}\~ao horizontal) com a holonomia (que mede o efeito global acumulado dessa n\~ao integrabilidade).

Seja \(\Omega\) a 2-forma de curvatura da conex\~ao \(\chi\), definida pela equa\c{c}\~ao de estrutura de Cartan:
\begin{equation}\label{eq:cartan-structure}
    \Omega = d\chi + \frac{1}{2}[\chi, \chi].
\end{equation}
Equivalentemente, \(\Omega(X,Y) = d\chi(X,Y) + [\chi(X), \chi(Y)]\) para campos vetoriais \(X,Y\) em \(P\).

\begin{teo}[Ambrose--Singer, 1953]\label{teo:ambrose-singer}
Seja \(P(M,G)\) um fibrado principal com conex\~ao \(\chi\) e curvatura \(\Omega\). Para todo \(u \in P\), a \'algebra de Lie \(\mathfrak{hol}_u(\chi)\) do grupo de holonomia \(\operatorname{Hol}_u(\chi)\) \'e o subespa\c{c}o de \(\mathfrak{g}\) gerado por todos os elementos da forma
\begin{equation}\label{eq:as-generators}
    \Omega_v(X,Y) \in \mathfrak{g},
\end{equation}
onde \(v \in P(u)\) pertence ao fibrado de holonomia de \(u\) e \(X,Y \in T_v P\) s\~ao vetores horizontais arbitr\'arios.
\end{teo}

Em termos mais operacionais: a \'algebra de Lie da holonomia \'e gerada por todos os valores que a curvatura \(\Omega\) assume sobre pares de vetores horizontais em pontos do fibrado de holonomia. Isto implica que:

\[
\boxed{\text{Curvatura } \Omega \;\; \text{``gera''} \;\; \text{Holonomia } \operatorname{Hol}(\chi)}
\]

\subsection{Esbo\c{c}o da Demonstra\c{c}\~ao}

A demonstra\c{c}\~ao completa encontra-se em Ambrose \& Singer (1953) e Kobayashi--Nomizu (1963, Cap. II, Teorema 8.1). Apresentamos aqui um esbo\c{c}o estruturado em quatro passos.

\begin{enumerate}
    \item \textbf{Fibrado de Holonomia.} Constr\'oi-se o subconjunto \(P(u) \subset P\) dos pontos que podem ser alcan\c{c}ados a partir de \(u\) por curvas horizontais. Demonstra-se que \(P(u)\) \'e um fibrado principal reduzido sobre \(M\) com grupo estrutural \(\operatorname{Hol}_u(\chi)\). A conex\~ao \(\chi\) restringe-se a uma conex\~ao em \(P(u)\).

    \item \textbf{Teorema de Redu\c{c}\~ao.} Pelo teorema de redu\c{c}\~ao de Cartan, uma conex\~ao em \(P\) reduz-se a um subfibrado \(Q \subset P\) se, e somente se, sua curvatura \(\Omega\) toma valores na \'algebra de Lie do grupo estrutural de \(Q\) quando avaliada em vetores tangentes a \(Q\). Aplicado a \(P(u)\), isto implica que \(\Omega_v(X,Y) \in \mathfrak{hol}_u(\chi)\) para todos \(v \in P(u)\) e \(X,Y\) horizontais.

    \item \textbf{A \'Algebra de Lie da Holonomia.} Seja \(\mathfrak{m}_u \subset \mathfrak{g}\) o subespa\c{c}o gerado por \(\{\Omega_v(X,Y) \mid v \in P(u),\; X,Y \text{ horizontais}\}\). Pelo passo anterior, \(\mathfrak{m}_u \subseteq \mathfrak{hol}_u(\chi)\). A inclus\~ao rec\'iproca \(\mathfrak{hol}_u(\chi) \subseteq \mathfrak{m}_u\) \'e a parte n\~ao trivial da demonstra\c{c}\~ao.

    \item \textbf{Inclus\~ao Rec\'iproca via a Equa\c{c}\~ao de Estrutura.} Considere a distribui\c{c}\~ao \(\Delta\) em \(P(u)\) gerada pelos campos horizontais e pelos campos fundamentais correspondentes a \(\mathfrak{m}_u\). Usando a equa\c{c}\~ao de estrutura de Cartan~\eqref{eq:cartan-structure} e a identidade de Bianchi \(D\Omega = 0\), demonstra-se que \(\Delta\) \'e involutiva e, pelo teorema de Frobenius, integr\'avel. A variedade integral m\'axima contendo \(u\) define uma redu\c{c}\~ao adicional do fibrado cujo grupo estrutural tem \'algebra de Lie \(\mathfrak{m}_u\). Como \(P(u)\) \'e a redu\c{c}\~ao m\'inima (j\'a que \(\operatorname{Hol}_u(\chi)\) \'e o menor grupo ao qual a conex\~ao se reduz), conclui-se que \(\mathfrak{hol}_u(\chi) \subseteq \mathfrak{m}_u\).
\end{enumerate}

\section{Consequ\^encia Financeira do Teorema de Ambrose--Singer}

\subsection{A Curvatura como Geradora de Arbitragem}

O Teorema de Ambrose--Singer fornece a base matem\'atica para a afirma\c{c}\~ao central da Teoria Geom\'etrica de Arbitragem:

\begin{teo}[Crit\'erio de Exist\^encia de Arbitragem]\label{teo:criterio-arbitragem}
Seja \(\chi\) uma conex\~ao de gauge sobre o fibrado de precifica\c{c}\~ao \(P(M,G)\) com curvatura \(\Omega\). Ent\~ao:
\begin{enumerate}
    \item Se \(\Omega \equiv 0\) (conex\~ao plana), ent\~ao \(\operatorname{Hol}(\chi)\) \'e trivial e \textbf{n\~ao existem} estrat\'egias de arbitragem.
    \item Se \(\Omega \not\equiv 0\), ent\~ao \(\operatorname{Hol}(\chi)\) \'e n\~ao trivial e \textbf{existe pelo menos uma} estrat\'egia de arbitragem com ganho n\~ao nulo.
\end{enumerate}
\end{teo}

\begin{proof}
(1) Se \(\Omega \equiv 0\), pelo Teorema de Ambrose--Singer a \'algebra de Lie \(\mathfrak{hol}(\chi) = \{0\}\), logo \(\operatorname{Hol}(\chi)\) \'e discreto. Como a holonomia restrita \'e a componente conexa da identidade, \(\operatorname{Hol}^0(\chi) = \{e\}\) e, para la\c{c}os contr\'ateis, o transporte paralelo \'e trivial.

(2) Se \(\Omega_v(X,Y) \neq 0\) para algum \(v \in P\) e \(X,Y\) horizontais, ent\~ao, pelo Teorema de Ambrose--Singer, este valor pertence a \(\mathfrak{hol}(\chi) \setminus \{0\}\). A exponencial \(\exp(t \,\Omega_v(X,Y)) \in \operatorname{Hol}^0(\chi)\) fornece uma fam\'ilia n\~ao trivial de elementos de holonomia, cada um correspondendo a uma estrat\'egia de arbitragem com ganho n\~ao nulo.
\end{proof}

\begin{prop}[Dimens\~ao e N\'umero de Estrat\'egias Independentes]
A dimens\~ao da \'algebra de Lie \(\mathfrak{hol}(\chi)\) (como espa\c{c}o vetorial real) \'e igual ao n\'umero m\'aximo de estrat\'egias de arbitragem linearmente independentes que o mercado admite:
\begin{equation}\label{eq:dim-arbitragem}
    \dim_{\mathbb{R}} \mathfrak{hol}(\chi) = N_{\text{arb}}.
\end{equation}
Em particular, \(\dim \mathfrak{hol}(\chi) \leq \dim \mathfrak{g}\).
\end{prop}

\subsection{O Fluxo Conceitual}

A cadeia de implica\c{c}\~oes que estrutura toda a teoria pode ser resumida no seguinte diagrama:

\[
\begin{array}{ccccc}
    \Omega \neq 0 & \xrightarrow{\text{Ambrose--Singer}} & \mathfrak{hol}(\chi) \neq \{0\} & \xrightarrow{\text{Exponencial}} & \operatorname{Hol}(\chi) \neq \{e\} \\[4pt]
    \text{Curvatura} & & \text{Geradores infinitesimais} & & \text{Holonomia finita} \\[4pt]
    & & \downarrow & & \downarrow \\[4pt]
    & & \multicolumn{3}{c}{\text{Estrat\'egias de Arbitragem}}
\end{array}
\]

Em palavras: \textbf{a curvatura ``gera'' a holonomia, e a holonomia ``parametriza'' as estrat\'egias de arbitragem.}

\section{Classifica\c{c}\~ao Alg\'ebrica de Estrat\'egias de Arbitragem}

\subsection{Estrutura da \'Algebra de Lie \(\mathfrak{hol}(\chi)\)}

A \'algebra de Lie \(\mathfrak{hol}(\chi)\) herda a estrutura de \(\mathfrak{g}\), mas \'e em geral um subconjunto pr\'oprio. Sua classifica\c{c}\~ao fornece informa\c{c}\~ao financeira precisa:

\begin{defi}[Dire\c{c}\~oes de Arbitragem]
Cada elemento n\~ao nulo \(A \in \mathfrak{hol}(\chi)\) define uma \textbf{dire\c{c}\~ao de arbitragem}. A estrat\'egia associada \'e obtida pela exponencia\c{c}\~ao \(\gamma_A(t) = \exp(tA)\), que descreve a evolu\c{c}\~ao do portf\'olio ao longo de um ciclo infinitesimal na dire\c{c}\~ao especificada por \(A\).
\end{defi}

\begin{prop}[Decomposi\c{c}\~ao Espectral das Estrat\'egias]
Seja \(\{T_1, \dots, T_k\}\) uma base da \'algebra de Lie \(\mathfrak{hol}(\chi)\), com \(k = \dim \mathfrak{hol}(\chi)\). Ent\~ao toda estrat\'egia de arbitragem infinitesimal pode ser expressa como combina\c{c}\~ao linear
\begin{equation}\label{eq:linear-comb}
    A = \sum_{i=1}^{k} \alpha_i \, T_i, \qquad \alpha_i \in \mathbb{R},
\end{equation}
e cada coeficiente \(\alpha_i\) representa a \textbf{aloca\c{c}\~ao} na \(i\)-\'esima estrat\'egia de arbitragem independente.
\end{prop}

\subsection{Classifica\c{c}\~ao por Estrutura da \'Algebra}

A natureza da \'algebra de Lie \(\mathfrak{hol}(\chi)\) determina propriedades qualitativas do conjunto de estrat\'egias de arbitragem dispon\'iveis:

\begin{center}
\begin{tabular}{@{}lll@{}}
\toprule
\textbf{Propriedade de \(\mathfrak{hol}(\chi)\)} & \textbf{Classificação} & \textbf{Interpretação Financeira} \\
\midrule
\(\mathfrak{hol} = \{0\}\)                    & Mercado sem arbitragem       & Todos os preços são justos; vale a Lei do Preço Único \\
\(\mathfrak{hol}\) abeliana                    & Arbitragem comutativa        & Estratégias independentes; ordem de execução irrelevante \\
\(\mathfrak{hol}\) nilpotente                  & Arbitragem hierárquica       & Ganhos decrescentes com repetição; saturação do mercado \\
\(\mathfrak{hol}\) simples                     & Arbitragem indecomponível    & Não há subestratégias isoláveis; risco sistêmico \\
\(\mathfrak{hol}\) semissimples                & Arbitragem decomposta        & Carteira de estratégias ortogonais entre si \\
\(\dim \mathfrak{hol} = \dim \mathfrak{g}\)    & Mercado saturado             & Número máximo de direções de arbitragem exploráveis \\
\bottomrule
\end{tabular}
\end{center}

\subsection{Invariantes Alg\'ebricos e Significado Financeiro}

\begin{prop}[Posto e N\'umero de Estrat\'egias]
O \textbf{posto} da \'algebra de Lie \(\operatorname{rank}(\mathfrak{hol}(\chi))\) (dimens\~ao m\'axima de uma sub\'algebra de Cartan) corresponde ao n\'umero m\'aximo de estrat\'egias de arbitragem \textbf{mutuamente comutativas}, i.e., que podem ser executadas simultaneamente sem interfer\^encia m\'utua.
\end{prop}

\begin{prop}[\'Indice de Arbitragem]
O \textbf{\'indice} \(\iota(\mathfrak{hol}) = \dim \mathfrak{hol} - \operatorname{rank}(\mathfrak{hol})\) mede o grau de n\~ao comutatividade das estrat\'egias de arbitragem. Um \'indice elevado indica que a ordem de execu\c{c}\~ao das estrat\'egias \'e relevante para o resultado financeiro final.
\end{prop}

\section{Exemplo Guiado: Mercado com Pre\c{c}os de It\^o}

\subsection{Modelagem do Mercado}

Considere um mercado com \(n\) ativos de risco cujos pre\c{c}os seguem processos de It\^o sob a medida f\'isica \(\mathbb{P}\):
\begin{equation}\label{eq:ito-prices}
    \frac{dS_t^{\,j}}{S_t^{\,j}} = \mu_t^{\,j} \, dt + \sum_{\alpha=1}^{d} \sigma_t^{\,j\alpha} \, dW_t^{\,\alpha}, \qquad j = 1, \dots, n,
\end{equation}
onde \(W_t = (W_t^1, \dots, W_t^d)\) \'e um movimento browniano \(d\)-dimensional, \(\mu_t = (\mu_t^1, \dots, \mu_t^n)\) \'e o vetor de retornos esperados (drifts), e \(\sigma_t = (\sigma_t^{\,j\alpha})\) \'e a matriz de volatilidades \(n \times d\).

A taxa livre de risco \'e denotada por \(r_t\). O grupo de gauge natural \'e \(G = \mathbb{R}^n\) (transla\c{c}\~oes nos log-pre\c{c}os), e a variedade base \(M\) \'e o espa\c{c}o de configura\c{c}\~oes dos pre\c{c}os.

\subsection{Conex\~ao de Precifica\c{c}\~ao}

A conex\~ao de gauge \(\chi\) \'e definida como a 1-forma que ``corrige'' o drift observado \(\mu_t\) para o drift neutro ao risco \(r_t \mathbf{1}\):
\begin{equation}\label{eq:pricing-connection}
    \chi_t = (\mu_t - r_t \mathbf{1}) \, dt.
\end{equation}
Aqui, \(\mathbf{1} = (1, \dots, 1)^\top \in \mathbb{R}^n\) e \(\chi_t\) toma valores na \'algebra de Lie \(\mathfrak{g} \cong \mathbb{R}^n\).

A condi\c{c}\~ao de aus\^encia de arbitragem (na formula\c{c}\~ao geom\'etrica) equivale a \(\chi\) ser uma conex\~ao plana, i.e., curvatura nula.

\subsection{Curvatura da Conex\~ao de Precifica\c{c}\~ao}

Como o grupo \(G = \mathbb{R}^n\) \'e abeliano, o colchete de Lie \'e identicamente nulo, e a equa\c{c}\~ao de estrutura de Cartan~\eqref{eq:cartan-structure} simplifica-se para:
\begin{equation}\label{eq:curvature-abelian}
    \Omega = d\chi.
\end{equation}

Em coordenadas locais, a curvatura \'e a 2-forma diferencial de \(\chi\). No contexto estoc\'astico, devemos considerar a formula\c{c}\~ao de Stratonovich para que o c\'alculo diferencial exterior seja v\'alido. A convers\~ao de It\^o para Stratonovich introduz um termo de corre\c{c}\~ao:
\begin{equation}\label{eq:ito-strat}
    \chi_t^{\text{Strat}} = \chi_t^{\text{It\^o}} - \frac{1}{2} d\langle \chi, W \rangle_t.
\end{equation}

Para o nosso caso, com \(\chi_t = (\mu_t - r_t \mathbf{1})\,dt\), n\~ao h\'a termo estoc\'astico em \(\chi\) (a conex\~ao \'e puramente um termo de drift), portanto as formula\c{c}\~oes de It\^o e Stratonovich coincidem para \(\chi\).

\subsection{C\'alculo Expl\'icito da Curvatura}

A curvatura \(\Omega = d\chi\) \'e obtida pela diferencial exterior da 1-forma \(\chi\). Em um contexto mais geral, se permitirmos que \(\mu_t\) e \(r_t\) dependam dos pre\c{c}os (e n\~ao apenas do tempo), a curvatura torna-se:
\begin{equation}\label{eq:curvature-market}
    \Omega_t = \sum_{j=1}^{n} \frac{\partial}{\partial S^j} (\mu_t - r_t \mathbf{1}) \wedge dS_t^{\,j}.
\end{equation}

Para o caso simplificado em que \(\mu_t\) e \(r_t\) s\~ao fun\c{c}\~oes apenas do tempo (determin\'isticos), temos:
\begin{equation}\label{eq:curvature-deterministic}
    \Omega_t = 0 \quad \Longleftrightarrow \quad \frac{\partial \mu_t^{\,j}}{\partial S^k} = 0 \;\; \forall j,k.
\end{equation}

\subsection{Interpreta\c{c}\~ao Financeira da Curvatura}

\begin{prop}[Condi\c{c}\~ao de Mercado sem Arbitragem Geom\'etrica]\label{prop:no-arbitrage-condition}
Em um mercado modelado pela Equa\c{c}\~ao~\eqref{eq:ito-prices} com conex\~ao de precifica\c{c}\~ao~\eqref{eq:pricing-connection}, as seguintes condi\c{c}\~oes s\~ao equivalentes:
\begin{enumerate}
    \item \(\Omega \equiv 0\) (curvatura nula).
    \item Existe um processo \(\lambda_t\) (pre\c{c}o de mercado do risco) tal que \(\mu_t^{\,j} - r_t = \sum_{\alpha=1}^{d} \sigma_t^{\,j\alpha} \lambda_t^{\,\alpha}\) para todo \(j\).
    \item Existe uma medida martingale equivalente \(\mathbb{Q}\) sob a qual os pre\c{c}os descontados s\~ao martingales.
\end{enumerate}
\end{prop}

\begin{proof}
A equival\^encia (2) \(\Leftrightarrow\) (3) \'e o Primeiro Teorema Fundamental de Precifica\c{c}\~ao de Ativos (FTAP). A equival\^encia (1) \(\Leftrightarrow\) (2) decorre do Teorema de Ambrose--Singer: \(\Omega = 0\) se, e somente se, \(\mathfrak{hol}(\chi) = \{0\}\), o que ocorre exatamente quando existe uma transforma\c{c}\~ao de gauge (mudan\c{c}a de medida) que anula a conex\~ao \(\chi\). No contexto abeliano, isto \'e equivalente \`a exist\^encia de \(\lambda_t\) satisfazendo a condi\c{c}\~ao de aus\^encia de arbitragem da Equa\c{c}\~ao~\eqref{eq:ito-prices}.
\end{proof}

\begin{prop}[Curvatura como Medida de Inefici\^encia]\label{prop:curvature-inefficiency}
A magnitude da curvatura \(\|\Omega_t\|\) (na norma de Hilbert--Schmidt) fornece uma \textbf{medida local de inefici\^encia do mercado}: quanto maior \(\|\Omega_t\|\), maior a intensidade das oportunidades de arbitragem dispon\'iveis na vizinhan\c{c}a do estado de mercado atual.
\end{prop}

\subsection{Exemplo Num\'erico Simples}

Considere um mercado com dois ativos (\(n=2\)) e um \'unico fator de risco (\(d=1\)), com par\^ametros constantes:
\[
\mu = \begin{pmatrix} 0{,}10 \\ 0{,}07 \end{pmatrix}, \quad
\sigma = \begin{pmatrix} 0{,}20 \\ 0{,}15 \end{pmatrix}, \quad
r = 0{,}05.
\]

O vetor de excesso de retorno \'e:
\[
\mu - r\mathbf{1} = \begin{pmatrix} 0{,}05 \\ 0{,}02 \end{pmatrix}.
\]

\begin{itemize}
    \item \textbf{Caso 1: Mercado completo e sem arbitragem.} Se existe \(\lambda\) tal que \(\mu - r\mathbf{1} = \sigma \lambda\), ent\~ao \(\lambda = \frac{0{,}05}{0{,}20} = \frac{0{,}02}{0{,}15} = 0{,}25\)? N\~ao: \(\frac{0{,}05}{0{,}20} = 0{,}25\) e \(\frac{0{,}02}{0{,}15} \approx 0{,}133\). Os pre\c{c}os de mercado do risco diferem -- logo h\'a oportunidade de arbitragem.

    \item \textbf{Curvatura n\~ao nula.} A conex\~ao \(\chi = (\mu - r\mathbf{1})dt\) tem curvatura \(\Omega = d\chi\). Se \(\mu\) \'e constante, \(d\mu = 0\) e \(\Omega = 0\). No entanto, se \(\mu = \mu(S)\) depende dos pre\c{c}os, a curvatura \'e n\~ao nula e a holonomia \'e n\~ao trivial.

    \item \textbf{Estrat\'egia de arbitragem.} Comprar o ativo com maior \'indice de Sharpe \(\frac{\mu^j - r}{\sigma^j}\) e vender o de menor \'indice, em propor\c{c}\~oes que anulem o risco total. O ganho \'e codificado pela holonomia ao longo do ciclo de trading.
\end{itemize}

\section{Resumo: Dicion\'ario Geometria \(\leftrightarrow\) Finan\c{c}as}

A tabela a seguir sintetiza as correspond\^encias fundamentais desenvolvidas neste cap\'itulo:

\begin{center}
\begin{tabular}{@{}>{\raggedright\arraybackslash}p{5.5cm}>{\raggedright\arraybackslash}p{5.5cm}@{}}
\toprule
\textbf{Conceito Geom\'etrico} & \textbf{Conceito Financeiro} \\
\midrule
Fibrado principal \(P(M,G)\) & Universo de todas as configurações de carteira possíveis \\
\addlinespace
Variedade base \(M\) & Espaço de estados do mercado (preços, volatilidades) \\
\addlinespace
Grupo estrutural \(G\) & Grupo de transformações de gauge (mudanças de numéraire, medidas) \\
\addlinespace
Conexão \(\chi\) & Estrutura de precificação (drifts ajustados ao risco) \\
\addlinespace
Curvatura \(\Omega = d\chi + \frac{1}{2}[\chi,\chi]\) & Intensidade local das oportunidades de arbitragem \\
\addlinespace
Transporte paralelo horizontal & Evolução ``neutra ao risco'' de uma carteira \\
\addlinespace
Laço \(\gamma\) em \(M\) & Ciclo fechado de trading \\
\addlinespace
Holonomia \(g_\gamma \in \operatorname{Hol}(\chi)\) & Ganho (ou perda) líquido ao completar o ciclo \\
\addlinespace
Grupo de holonomia \(\operatorname{Hol}(\chi)\) & Conjunto de todos os resultados financeiros possíveis de ciclos de trading \\
\addlinespace
Holonomia restrita \(\operatorname{Hol}^0(\chi)\) & Resultados de ciclos de trading continuamente deformáveis entre si \\
\addlinespace
Álgebra de Lie \(\mathfrak{hol}(\chi)\) & Geradores infinitesimais das estratégias de arbitragem \\
\addlinespace
Base de \(\mathfrak{hol}(\chi)\) & Conjunto maximal de estratégias de arbitragem independentes \\
\addlinespace
\(\dim \mathfrak{hol}(\chi)\) & Número de estratégias de arbitragem linearmente independentes \\
\addlinespace
Teorema de Ambrose--Singer & A curvatura gera todas as estratégias de arbitragem (infinitesimais) \\
\addlinespace
\(\Omega \equiv 0\) (conexão plana) & Ausência de arbitragem (vale o FTAP) \\
\addlinespace
\(\Omega \not\equiv 0\) & Existência de pelo menos uma estratégia de arbitragem não trivial \\
\addlinespace
Equação de estrutura de Cartan & Relação entre precificação (drift) e arbitragem (curvatura) \\
\bottomrule
\end{tabular}
\end{center}

\subsection*{Refer\^encias do Cap\'itulo}

\begin{itemize}
    \item \textsc{Ambrose, W.; Singer, I. M.} A theorem on holonomy. \emph{Transactions of the American Mathematical Society}, v.~75, n.~3, p.~428--443, 1953. DOI: \texttt{10.2307/1990721}.
    \item \textsc{Kobayashi, S.; Nomizu, K.} \emph{Foundations of Differential Geometry}. Volumes I (1963) e II (1969). Interscience Publishers (John Wiley \& Sons), New York.
    \item \textsc{Besse, A. L.} \emph{Einstein Manifolds}. Springer-Verlag, 1987. (Cap\'itulo 10: Holonomy Groups.)
    \item \textsc{Joyce, D. D.} \emph{Riemannian Holonomy Groups and Calibrated Geometry}. Oxford University Press, 2007.
    \item \textsc{Delbaen, F.; Schachermayer, W.} \emph{The Mathematics of Arbitrage}. Springer Finance, 2006.
    \item \textsc{Ilinski, K.} \emph{Physics of Finance: Gauge Modelling in Non-Equilibrium Pricing}. John Wiley \& Sons, 2001.
    \item \textsc{Young, K.} Foreign exchange market as a lattice gauge theory. \emph{American Journal of Physics}, v.~67, n.~10, p.~862--868, 1999.
\end{itemize}


